\par\medskip
\noindent\textbf{Observación 5.8.}
Sea $\mathfrak X$ un esquema formal $S$-ádico esencialmente de tipo
finito. Para un punto $x\in\mathfrak X$, consideremos también el
esquema formal $S$-ádico
\[
\mathfrak X_{(\widehat{\sisi{X}{x}})}
=\varinjlim_\nu
 \operatorname{Spec}\bigl(\widehat{\mathcal O_{X_\nu,x}}\bigr).
\]
Hay morfismos naturales
\[
\mathfrak X_{(\widehat{\sisi{X}{x}})}
\longrightarrow\mathfrak X_{(x)}
\longrightarrow\mathfrak X.
\]
En general, $\widehat{\mathcal O_{X_0,x}}$ no es esencialmente de
tipo finito sobre $S_{(y)}$, donde $y$ es la imagen de $x$ por
$\mathfrak X\longrightarrow S$.

Para simplificar, pongamos
\[
A_\nu=\mathcal O_{X_\nu,x},
\qquad R_\nu=\mathcal O_{S_\nu,y}.
\]
El morfismo natural
\[
\Omega_{A_\nu|R_\nu}
\longrightarrow\Omega_{\widehat A_\nu|R_\nu}
\]
induce un isomorfismo
\[
\widehat{\Omega_{A_\nu|R_\nu}}
\xrightarrow{\ \sim\ }
\sisi{\widehat{\widehat\Omega_{A_\nu|R_\nu}}}{\widehat{\Omega_{\widehat A_\nu|R_\nu}}},
\]
y por tanto
\[
\varprojlim_\nu\widehat{\Omega_{A_\nu|R_\nu}}
\xrightarrow{\ \sim\ }
\varprojlim_\nu
 \sisi{\widehat{\Omega_{\widehat A_\nu|R_\nu}}}{\widehat{\Omega_{\widehat A_\nu|R_\nu}}}.
\]
Aquí el sombrero exterior denota la completación separada de un
módulo para la topología inducida por el anillo. Por consiguiente, si
definimos
\[
\widehat\Omega_{\mathfrak X_{(\widehat{x})}|S}
=\varprojlim_\nu
 \sisi{\widehat{\Omega_{
   \widehat{\mathcal O}_{X_\nu,x}|\mathcal O_{S_\nu,y}}}}{\widehat{\Omega_{
   \widehat{\mathcal O}_{X_\nu,x}|\mathcal O_{S_\nu,y}}}},
\]
entonces $\widehat\Omega_{\mathfrak X_{(\widehat{x})}|S}$ es un
$\mathfrak X_{(\widehat{x})}$-módulo coherente.

%%%%%%%%%%  scan idx 133  |  folio 122  %%%%%%%%%%

Definamos
$J_{\mathfrak X_{(\widehat{x})}}^p
 (\mathfrak X_{(\widehat{x})}|S)$
como el subesquema formal cerrado de
$\mathfrak X_{(\widehat{x})}$ definido por
\[
\underline I_{\mathfrak X_{(\widehat{x})}}^p
 (\mathfrak X_{(\widehat{x})}|S)
=\underline I_{\mathfrak X_{(\widehat{x})}}^p
 (\widehat\Omega_{\mathfrak X_{(\widehat{x})}|S}).
\]
Entonces
\[
J_{\mathfrak X_{(x)}}^p(\mathfrak X_{(x)}|S)_{(\widehat{x})}
\simeq
J_{\mathfrak X_{(\widehat{x})}}^p
 (\mathfrak X_{(\widehat{x})}|S).
\]
Si $x$ es un punto aislado de
$J_{\mathfrak X}^p(\mathfrak X|S)$, el morfismo natural
\[
J_{\mathfrak X_{(\widehat{x})}}^p
 (\mathfrak X_{(\widehat{x})}|S)
\longrightarrow
J_{\mathfrak X_{(x)}}^p(\mathfrak X_{(x)}|S)
\]
es un isomorfismo.

Para formar la ``imagen'' de $J_{\mathfrak X}^p(\mathfrak X|S)$ por
$f:\mathfrak X\longrightarrow S$, podríamos tomar simplemente el
subesquema formal de $S$ definido por el haz de ideales
\[
\ker\left(
 \mathcal O_S\longrightarrow
 \mathcal O_{\mathfrak X}/
 \underline I_{\mathfrak X}^p(\mathfrak X|S)
\right).
\]
Para nuestros fines, sin embargo, es más conveniente volver a utilizar
el ideal de Fitting.
